\section{Discusión}

\subsection{Significado físico y cumplimiento de los objetivos}
Los cinco máximos de la Tabla~III corresponden a frecuencias
en las que la columna de aire entra en resonancia con la
onda sonora generada por el teléfono: la onda emitida se superpone con sus
reflexiones en el fondo y en la abertura, y se forma una onda
estacionaria de gran amplitud. Los cocientes $f_n/f_1$
resultaron $1$, $2.98$, $4.81$, $7.15$ y $9.02$, es
decir, aproximadamente $1, 3, 5, 7$ y $9$. Solo aparecen
armónicos impares porque las condiciones de frontera del tubo
cerrado (nodo de desplazamiento en el fondo y antinodo en la
abertura) obligan a que en $L_f$ quepa un número impar de
cuartos de longitud de onda, $L_f = n\lambda_n/4$. Los
armónicos pares exigirían nodos o antinodos en ambos extremos,
lo que es incompatible con esta geometría. Así se cumplió el
objetivo de generar, detectar y asociar las resonancias con los
armónicos permitidos.

El ajuste $f_n = mn + b$ (Tabla~IV) respalda este modelo:
$R^2 = 0.9987$ y un intercepto
$b = (-13 \pm 28)\,\mathrm{Hz}$ compatible con cero dentro de
su incertidumbre, como exige $f_n = nf_1$. La pendiente
$f_1 = (227.5 \pm 4.8)\,\mathrm{Hz}$ dio
$v = (341.3 \pm 7.7)\,\mathrm{m/s}$. La diferencia con
$v_T = (345.9 \pm 0.6)\,\mathrm{m/s}$ es de
$4.6\,\mathrm{m/s}$ ($1.34\,\%$), menor que la
incertidumbre $\Delta v$, por lo que ambos valores son
compatibles y se cumple el tercer objetivo.

En la Parte~II la superficie del agua actúa como extremo
cerrado y la frecuencia de resonancia aumenta al acortarse la
columna: de $218.23\,\mathrm{Hz}$ con
$L' = 37.50\,\mathrm{cm}$ a $596.73\,\mathrm{Hz}$ con
$L' = 14.58\,\mathrm{cm}$ (Tabla~VI). La frecuencia se
multiplicó por $2.73$ mientras que la longitud se redujo en
un factor $2.57$, tendencia cercana a la proporcionalidad
inversa de la Ec.~(9). Los niveles de los picos se mantuvieron
entre $64$ y $65\,\mathrm{dB}$ en los casos reportados, lo que
indica una excitación y detección comparables en todos los
niveles de agua.

\subsection{Discrepancias respecto al modelo teórico}
Con $L_f = 0.375\,\mathrm{m}$ y
$v_T = 345.875\,\mathrm{m/s}$, el modelo predice
$f_1 = v_T/(4L_f) = 230.6\,\mathrm{Hz}$. El pico medido
($225.22\,\mathrm{Hz}$) queda $2.3\,\%$ por debajo y la
pendiente del ajuste $1.3\,\%$ por debajo. Los cocientes
$f_n/n$ van de $216.8$ a $230.1\,\mathrm{Hz}$ (alrededor de
$224.4\,\mathrm{Hz}$), sin una tendencia monótona con $n$. El
menor corresponde a $n=5$ ($1083.83\,\mathrm{Hz}$). Esta
dispersión es mucho mayor que la resolución espectral de la
aplicación ($1.35\,\mathrm{Hz}$, Fig.~5), así que
probablemente se debe a picos anchos y a la respuesta en
frecuencia no plana del altavoz y el micrófono del teléfono,
que desplazan el máximo aparente. Además, la incertidumbre de
$v$ proviene casi toda de la frecuencia
($\Delta f/f = 2.1\,\%$) y muy poco de la geometría
($\Delta L_f/L_f = 0.14\,\%$).

En la Parte~II el producto $fL'$, que el modelo predice
constante e igual a $v_T/4 \approx 86.5\,\mathrm{Hz\cdot m}$,
varía de $81.8$ a $87.0\,\mathrm{Hz\cdot m}$ (de
$-5.4\,\%$ a $+0.6\,\%$). Las columnas largas, con
resonancias más bajas, se alejan más del modelo, mientras que
las dos más cortas coinciden con él dentro de $0.6\,\%$. Esto
sugiere un efecto más importante a baja frecuencia, como la
respuesta del altavoz del teléfono, o un efecto de la abertura
parcialmente cubierta por el teléfono. Con la misma columna de
aire (sin agua) se midió $225.22\,\mathrm{Hz}$ en la
Tabla~III y $218.23\,\mathrm{Hz}$ en la Tabla~VI, una
diferencia de $3.1\,\%$ con barridos distintos
(100--2300\,Hz en 20\,s frente a 200--1000\,Hz en 10\,s). Esto
indica que la reproducibilidad real de la ubicación del máximo
es peor que la incertidumbre estadística del ajuste. El ajuste
exponencial de la Fig.~7 ($R^2 = 0.9813$) describe de forma
empírica la tendencia en el intervalo medido; la ley física
esperada es $f = v/(4L')$.

La comparación con los diapasones es especialmente sensible.
Como $V = \tfrac{\pi D^2}{4}\left(L_f - \tfrac{v_T}{4f}\right)$
depende de la diferencia entre dos longitudes cercanas, un
error de $1\,\%$ en $v$ desplaza el volumen unos
$6.6\,\mathrm{cm^3}$ para $256\,\mathrm{Hz}$. Los volúmenes
teóricos de la Tabla~VII ($73.10\,\mathrm{cm^3}$ para
$256\,\mathrm{Hz}$ y $382.60\,\mathrm{cm^3}$ para
$480\,\mathrm{Hz}$) son cercanos a los niveles de $75.0$ y
$375.0\,\mathrm{cm^3}$ de la Tabla~VI. En esos niveles la
resonancia medida fue de $242.92$ y $468.13\,\mathrm{Hz}$,
es decir, $5.1\,\%$ y $2.5\,\%$ por debajo de la frecuencia
del diapasón. Por ello, en la práctica se requiere ajustar el
nivel del agua respecto del valor teórico.

\subsection{Limitaciones y fuentes de error}
Entre las fuentes sistemáticas están la corrección de extremo
$0.3D$, que es una aproximación para longitudes de onda
grandes frente al diámetro (para $n=9$,
$\lambda \approx 0.17\,\mathrm{m}$ es solo unas $3.3$ veces
$D$), el teléfono cubriendo parcialmente la abertura, la
respuesta no plana de sus transductores y el uso de $v_T$
lineal en $T$, que supone aire seco y temperatura uniforme
medida una sola vez ($\pm 1\,^{\circ}\mathrm{C}$). Entre las
aleatorias están el ruido ambiental del laboratorio, la lectura
manual del máximo en el espectro y la incertidumbre de
$\pm 2\,\mathrm{cm^3}$ al agregar el agua, que hace crecer
$\Delta L'$ de $0.05$ a $0.20\,\mathrm{cm}$. Además, muchas
mediciones se hicieron una sola vez, sin repeticiones que
permitan estimar una incertidumbre estadística.

Como mejoras, se recomiendan promediar varios barridos, usar
barridos lentos y de banda estrecha alrededor de cada $f_n$,
emplear altavoz y micrófono externos, medir el volumen con
material más preciso y registrar la temperatura durante todo el
experimento. También convendría ajustar $f$ frente a $1/L'$ con
más niveles de agua para contrastar la pendiente con $v_T/4$.